\chapter*{Prefacio}
\addcontentsline{toc}{chapter}{Prefacio}

Existe esta creencia, tanto dentro como fuera del gremio, de que aquellas
personas que nos dedicamos a las matem\'aticas no hacemos otra cosa mas que
pensar en ellas. Si bien no concuerdo del todo con esta creencia (y no lo digo
solo por ser un \textit{casi matem\'atico}), algo que no puedo negar es que, por
m\'as rid\'iculo que parezca, trataremos de involucrar nuestra \'area de estudio
en casi cualquier situaci\'on de nuestras vidas.

Guiados quiz\'as por nuestra inagotable curiosidad o simplemente por la
necesidad de un reto en nuestro d\'ia, es impresionante lo tan capaces que somos
para plantearnos preguntas que, si alguien en alg\'un momento se habr\'a
cuestionado, dudo que \'este encuentre divertido invertirle horas y horas solo
por la satisfacci\'on de decir \textit{``listo, no pod\'ia quedarme con la
duda''}.

En mi experiencia, la mina de oro de donde surgen est\'as preguntas pueden ser
situaciones tan simples como una reuni\'on de amigos para disfrutar de un juego
de mesa. No es casualidad que esta actividad que disfruto tanto dispare mi
creatividad con interrogantes respecto a estrategias y partidas \'optimas (y, de
vez en cuando, debates sobre las injusticias de la vida). En mi opini\'on, es
aqu\'i donde la combinatoria sale a relucir como una rama de las matem\'aticas
tan sencilla de comprender mediante ejemplos, pero tan compleja de domar en su
totalidad.

En ese esp\'iritu, esta tesis trata precisamente de c\'omo un juego tradicional,
como lo es polic\'ias y ladrones, puede recibir el tratamiento m\'etodico y
riguroso de las matem\'aticas en busca un conocimiento m\'as profundo del mismo.
Las respuestas descubiertas en este trabajo pretenden esclarecer los secretos
detr\'as de un juego universal, de modo que, al estilo de un mago en su truco de
magia, podemos retar a aquellas personas no tan versadas en el tema a partidas
en las que, gracias al dios de las matem\'aticas, siempre ganaremos. Al final
del d\'ia, si no quieres quedarte con la duda de si te est\'an estafando o no,
puede que sea pertinente que comiences a pensar como matem\'atico.

As\'i pues, el texto inicialmente revisa las bases del juego y su funcionamiento
a trav\'es de la teor\'ia de gr\'aficas, de las cuales se obtienen los
resultados cl\'asicos del \'area. En busca de expandir dichos resultados,
introducimos el juego sobre las gr\'aficas de fichas como una variante por
equipos, de manera que podemos deducir par\'ametros de gr\'aficas grandes en
funci\'on de propiedades de gr\'aficas peque\~nas. De nuestro trabajo en
gr\'aficas de fichas de \'arboles podemos dar cotas para el n\'umero policiaco
de una partida de polic\'ias y ladrones sobre esa familia de gr\'aficas.
Adem\'as, obtenemos algoritmos f\'aciles que nos ayudar\'an, a manera de
exploraci\'on, a determinar con exactitud el n\'umero de polic\'ia de una
gr\'afica de fichas de un \'arbol. Finalmente, agregamos, a manera exploraci\'on
en una gr\'afica, una cota junto con una estrategia para el n\'umero policiaco
de la gr\'afica de $2$ fichas de Petersen.